% ===================================================================
%  TABLEAU N° 3 — L'enfance et la jeunesse.
%
%  Traduction, faite en 2026, en francais standard du Quebec, sur
%  l'ido de texto/io. Regles de la colonne : voir 10-tableau-01.tex.
%
%  DEUX SCENES, DEUX NUMEROTATIONS. Comme le tableau 4, celui-ci
%  reunit deux scenes de famille — « Premiere scene » (un bapteme au
%  village) et « Deuxieme scene » (une fete publique) — dont chacune
%  recommence sa numerotation a 1. Les cles portent donc c1 et c2 :
%  t03-c1-01-1 et t03-c2-01-1 sont deux alineas « 1 » differents.
%
%  LES SIX CHOIX PROPRES A CE TABLEAU :
%    (6)  chaussettes ......... bas
%    (2)  tricot .............. chandail
%    (10) veste ............... veston
%    (52) sous ................ pièces de monnaie
%    (46) tir ................. stand de tir
%    (71) vieille femme ....... vieille dame
%
%  LE PIEGE EST DANS (10) ET (17), ET IL SE RETOURNE. Le fac-simile
%  francais dit « veste » (10) pour le vetement plie sur la berge, et
%  « gilet » (17) pour celui que le retardataire tient a la main. Au
%  Quebec les deux mots ont GLISSE D'UN CRAN : la « veste » y est le
%  vetement sans manches — ce que la France appelle gilet — et le
%  vetement a manches se dit « veston ». Traduire « veste » par
%  « veste » aurait donc designe le mauvais objet, et fait de (10) et
%  (17) le meme vetement. D'ou : (10) veston, (17) gilet.
%
%  « SOULIERS » RESTE. C'est l'inverse du piege precedent : le mot du
%  fac-simile est aussi celui du Quebec d'aujourd'hui, la ou la France
%  dirait plutot « chaussures ». On ne le remplace pas pour la seule
%  raison qu'il est ancien : il ne l'est pas ici.
%
%  LES CHOSES DE 1926 NE BOUGENT PAS : le corset, la redingote, le
%  chapeau haut-de-forme, la soutane, le presbytere, les chevaux de
%  bois et le mat de cocagne sont ce que la gravure montre. On ne
%  modernise que le mot, jamais l'objet.
% ===================================================================

%%K t03-apar-1 apar
\VUsaut{3.0mm}
\VUtitre{15.0pt}{60}{TABLEAU N\textsuperscript{o} 3}
\VUsaut{1.6mm}
\VUtitre{11.4pt}{0}{L'enfance et la jeunesse.}
\VUsaut{3.4mm}

%%K t03-apar-2 apar
(Les tableaux 3 et 4 sont combinés de telle sorte qu'ils présentent, en quatre
\VUgras{scènes de famille}, les notions essentielles sur le \VUgras{temps
qu'il fait}, l'\VUgras{âge}, la \VUgras{famille}, le \VUgras{vêtement} et la
\VUgras{condition}.)

\VUsaut{4.0mm}
%%K t03-tit-1 sub
\VUcentre{12.0pt}{0}{\textit{Première scène.}}
\VUsaut{3.0mm}

%%K t03-tit-2 sub
\VUpk{10.2pt}{40}{Un baptême au village.}
\VUsaut{2.4mm}

%%K t03-tit-3 sub
\VUpk{10.2pt}{40}{Le printemps.}
\VUsaut{3.0mm}

%%K t03-c1-01-1 p
1. --- Nous sommes au \VUgras{printemps}, dans les mois de \VUgras{mars},
d'\VUgras{avril} ou de \VUgras{mai}, un des jours de la \VUgras{semaine} :
\VUgras{lundi}, \VUgras{mardi} ou \VUgras{mercredi}. Il est \VUgras{dix
heures} du matin.

%%K t03-c1-02-1 p
2. --- Il fait beau \VUgras{temps}, l'\VUgras{air} est pur. Le soleil brille.
Quelques \VUgras{petits nuages} \textsuperscript{(2)} montent dans le
\VUgras{ciel} \textsuperscript{(1)} bleu.

%%K t03-c1-03-1 p
3. --- Les \VUgras{arbres} ont à peine des \VUgras{feuilles}. Les minces
\VUgras{peupliers} \textsuperscript{(3)} et le grand \VUgras{platane}
\textsuperscript{(17)} n'ont encore que des bourgeons.

%%K t03-c1-04-1 p
4. --- Les \VUgras{hirondelles} \textsuperscript{(4)} sont de retour et
voltigent en pépiant joyeusement autour de la \VUgras{flèche}
\textsuperscript{(7)} du \VUgras{clocher} \textsuperscript{(5)} qui couronne
l'\VUgras{église} \textsuperscript{(6)} du village.

%%K t03-tit-4 sub
\VUsaut{3.4mm}
\VUpk{10.2pt}{40}{L'église.}
\VUsaut{3.0mm}

%%K t03-c1-05-1 p
5. --- Elle est bien simple, cette église, avec son grand \VUgras{portail}
\textsuperscript{(13)} et sa grosse \VUgras{horloge} \textsuperscript{(8)}. La
\VUgras{petite aiguille} \textsuperscript{(9)} est sur le chiffre X : elle
marque les \VUgras{heures}. La \VUgras{grande} \textsuperscript{(10)} est sur
XII (la mi-heure); elle marque les \VUgras{minutes}.

%%K t03-c1-06-1 p
6. --- Dans le clocher, les \VUgras{cloches} \textsuperscript{(11)} sonnent
gaiement. Au sommet du toit se dresse une petite \VUgras{croix}
\textsuperscript{(12)} de pierre.

%%K t03-c1-07-1 p
7. --- L'église n'a pas seulement un grand \VUgras{portail}
\textsuperscript{(13)}, elle a aussi une petite porte de côté. C'est par cette
petite porte que monsieur le \VUgras{curé} \textsuperscript{(15)}, vêtu de sa
\VUgras{soutane}, sort de la \VUgras{sacristie} \textsuperscript{(14)} et se
rend au \VUgras{presbytère} \textsuperscript{(19)}.

%%K t03-c1-08-1 p
8. --- Près de lui, voici la \VUgras{laitière} \textsuperscript{(24)} avec son
\VUgras{âne} \textsuperscript{(23)}. Elle revient de la ville, où elle a
apporté du bon lait pour les enfants.

%%K t03-tit-5 sub
\VUsaut{4.0mm}
\VUpk{10.2pt}{40}{Le verger.}
\VUsaut{3.4mm}

%%K t03-c1-09-1 p
9. --- Monsieur Aliboron (l'âne) est tout content de cette course du matin. Il
respire avec délices l'air embaumé par le parfum des \VUgras{fleurs}
\textsuperscript{(20)} qui couvrent les \VUgras{arbres fruitiers}
\textsuperscript{(18)} du \VUgras{verger} voisin. Ils sont tout roses et tout
blancs, ces \VUgras{pêchers} \textsuperscript{(18)} et ces
\VUgras{abricotiers} \textsuperscript{(19)}, et ils font espérer une bonne
récolte.

%%K t03-c1-10-1 p
10. --- Pendant ce temps, les petites \VUgras{abeilles} \textsuperscript{(22)}
de la \VUgras{ruche} \textsuperscript{(21)} vont y butiner leur doux
\VUgras{miel} en bourdonnant.

%%K t03-c1-11-1 p
11. --- Mais que se passe-t-il? Voici les enfants du village qui accourent et
font fuir les \VUgras{oies} \textsuperscript{(25)} effrayées. C'est le cortège
qui sort de l'église après le \VUgras{baptême} de la fille du notaire.

%%K t03-c1-12-1 p
12. --- Le petit Jacques, dans son \VUgras{berceau} \textsuperscript{(26)}, se
réveille au joyeux son des cloches, et sa sœur Madeleine, qui veut voir elle
aussi le cortège du baptême, demande à sa mère de venir la coiffer et de
l'habiller sur le seuil de la maison.

%%K t03-c1-13-1 p
13. --- La mère y consent. Elle met son \VUgras{foulard}
\textsuperscript{(28)}, son \VUgras{corsage} \textsuperscript{(29)} et son
\VUgras{tablier} \textsuperscript{(30)}, et elle vient s'asseoir sur une petite
chaise devant la porte. Madeleine n'a que son \VUgras{jupon}
\textsuperscript{(31)} et son \VUgras{corset} \textsuperscript{(32)}.

%%K t03-c1-14-1 p
14. --- Comme elle est heureuse! Elle en oublie ses fleurs préférées, ses
\VUgras{œillets} \textsuperscript{(34)}, où se posent les \VUgras{papillons}
\textsuperscript{(35)}, ses \VUgras{tulipes} \textsuperscript{(36)}, ses
\VUgras{muguets} \textsuperscript{(37)} dans leurs \VUgras{pots}
\textsuperscript{(38)}, sur le \VUgras{mur} \textsuperscript{(33)}, et ses
\VUgras{roses} \textsuperscript{(39)}, qui forment une si belle haie. Elle
examine la riche toilette des dames du cortège.

%%K t03-c1-15-1 p
15. --- Les enfants du village ne font pas très attention aux vêtements. Ils se
précipitent et se bousculent les uns les autres pour avoir des \VUgras{bonbons}
\textsuperscript{(55)} ou des \VUgras{pièces de monnaie} \textsuperscript{(52)}.
L'un d'eux pleure \textsuperscript{(40)}.

%%K t03-c1-16-1 p
16. --- Un autre a marché sur la patte du \VUgras{chien} \textsuperscript{(42)},
qui se sauve en jappant et fait fuir la \VUgras{poule} \textsuperscript{(43)} et
ses \VUgras{poussins} \textsuperscript{(44)}.

%%K t03-tit-6 sub
\VUsaut{5.0mm}
\VUpk{10.2pt}{40}{Le cortège du baptême.}
\VUsaut{4.0mm}

%%K t03-c1-17-1 p
17. --- En tête du cortège marche la \VUgras{nourrice} \textsuperscript{(45)}.
Elle a une belle \VUgras{coiffe} \textsuperscript{(46)} à longs rubans. Elle
porte le \VUgras{nouveau-né} \textsuperscript{(47)}, tout habillé de
\VUgras{dentelles} \textsuperscript{(48)}.

%%K t03-c1-18-1 p
18. --- Derrière la nourrice viennent le \VUgras{parrain}
\textsuperscript{(49)} et la \VUgras{marraine} \textsuperscript{(53)}. Ils
distribuent les bonbons et les pièces de monnaie. Le parrain a des
\VUgras{gants} \textsuperscript{(51)} et un \VUgras{chapeau haut-de-forme}
\textsuperscript{(50)}. La marraine tient à la main un beau \VUgras{bouquet}
\textsuperscript{(54)}.

%%K t03-c1-19-1 p
19. --- Derrière eux, on aperçoit l'\VUgras{oncle} \textsuperscript{(56)} et la
\VUgras{tante} \textsuperscript{(58)} de l'enfant. La tante a un élégant
\VUgras{chapeau} \textsuperscript{(59)}, l'oncle une \VUgras{redingote}
\textsuperscript{(57)} noire et un gilet blanc. Voici ensuite le
\VUgras{cousin} \textsuperscript{(60)} et la \VUgras{cousine}
\textsuperscript{(61)}. Celle-ci tient par la main la \VUgras{sœur}
\textsuperscript{(62)} du nouveau-né, la petite Berthe.

%%K t03-c1-20-1 p
20. --- L'autre \VUgras{frère} \textsuperscript{(63)} de Berthe marche
derrière. Il donne la main à une \VUgras{parente} \textsuperscript{(64)}. Les
\VUgras{amis} \textsuperscript{(65)} de la famille viennent ensuite.

%%K t03-tit-7 sub
\VUsaut{4.6mm}
\VUpk{10.2pt}{40}{Les spectateurs.}
\VUsaut{3.4mm}

%%K t03-c1-21-1 p
21. --- Près du cortège, voici un \VUgras{mendiant} \textsuperscript{(66)}
appuyé sur sa \VUgras{béquille} \textsuperscript{(67)}. Il demande l'aumône (il
mendie).

%%K t03-c1-22-1 p
22. --- De l'autre côté, il y a un petit garçon très \VUgras{poli}
\textsuperscript{(68)}. Il salue et dit « \VUgras{bonjour} » aux personnes
qu'il connaît.

%%K t03-c1-23-1 p
23. --- Les villageois sont sortis de chez eux pour voir passer le cortège du
baptême. On voit un \VUgras{paysan} \textsuperscript{(69)} avec son
\VUgras{bonnet de coton} et, à côté, une \VUgras{paysanne}
\textsuperscript{(70)}. Puis une \VUgras{vieille dame} \textsuperscript{(71)}
appuyée sur son \VUgras{bâton} \textsuperscript{(72)}. Enfin le
\VUgras{meunier} \textsuperscript{(73)} avec son grand \VUgras{chapeau de
feutre} \textsuperscript{(74)} et une jeune paysanne chaussée de \VUgras{sabots}
\textsuperscript{(75)}, qui apprend à marcher à son tout-petit.

%%K t03-c1-24-1 p
24. --- Cette scène est très amusante.

\VUsaut{4.0mm}
%%K t03-tit-8 sub
\VUcentre{12.0pt}{0}{\textit{Deuxième scène.}}
\VUsaut{3.0mm}

%%K t03-tit-9 sub
\VUpk{10.2pt}{40}{Une fête publique.}
\VUsaut{2.4mm}

%%K t03-tit-10 sub
\VUpk{10.2pt}{40}{Jeux nautiques et régates.}
\VUsaut{3.0mm}

%%K t03-c2-01-1 p
1. --- L'\VUgras{été} est arrivé. Le soleil brûlant du mois de \VUgras{juin}
(juillet ou \VUgras{août}) donne aux jeunes l'envie de se baigner et de faire
du canot.

%%K t03-c2-02-1 p
2. --- Sur la rive droite du fleuve, les canotiers se déshabillent pour prendre
part aux jeux nautiques et aux régates.

%%K t03-c2-03-1 p
3. --- L'un d'eux enlève ses \VUgras{souliers} \textsuperscript{(1)}; il les
délace (il les déboutonne). Ses deux camarades sont déjà prêts. Ils n'ont plus
que leur \VUgras{chandail} \textsuperscript{(2)} et leur \VUgras{caleçon}
\textsuperscript{(3)}. Ils causent en attendant leurs amis. Ils parlent de la
foire et de la fête qui a lieu sur l'autre rive.

%%K t03-c2-04-1 p
4. --- Par terre, à côté d'eux, sont posés les \VUgras{vêtements} de leurs
camarades : une \VUgras{paire} de \VUgras{bottines} \textsuperscript{(4)}, des
\VUgras{souliers} \textsuperscript{(5)}, des \VUgras{bas}
\textsuperscript{(6)}, des chapeaux de \VUgras{feutre} \textsuperscript{(7)} et
de \VUgras{paille} \textsuperscript{(8)}, et une \VUgras{cravate}
\textsuperscript{(9)}.

%%K t03-c2-05-1 p
5. --- Un des jeunes gens a soigneusement plié son \VUgras{veston}
\textsuperscript{(10)}. Il le pose sur une serviette de toilette, dans laquelle
son voisin enveloppera bientôt les deux \VUgras{vêtements}.

%%K t03-c2-06-1 p
6. --- Un sixième garçon est en retard. Il a encore sa \VUgras{casquette}
\textsuperscript{(11)} sur la tête. Il n'a enlevé ni sa \VUgras{chemise}
\textsuperscript{(12)}, ni son \VUgras{col} \textsuperscript{(13)}, ni ses
\VUgras{manchettes} \textsuperscript{(14)}. Il tient son \VUgras{gilet}
\textsuperscript{(17)} à la main et il a encore son \VUgras{pantalon}
\textsuperscript{(16)} et ses \VUgras{bretelles} \textsuperscript{(15)}. Il ne
sera certainement pas prêt pour la prochaine course.

%%K t03-c2-07-1 p
7. --- Près des jeunes gens, un sentier bordé de \VUgras{chardons}
\textsuperscript{(18)} mène au bord de la rivière, où le père Jacques prépare,
dans les \VUgras{roseaux} \textsuperscript{(19)}, le \VUgras{feu d'artifice}
\textsuperscript{(20)} qu'on tirera le soir.

%%K t03-tit-11 sub
\VUsaut{4.4mm}
\VUpk{10.2pt}{40}{La course.}
\VUsaut{3.4mm}

%%K t03-c2-08-1 p
8. --- Sur la rivière, quelques dames, à l'abri de leur ombrelle, se promènent
en \VUgras{barque} \textsuperscript{(21)}. Elles suivent la course, qui est très
intéressante. Dans chaque \VUgras{yole} \textsuperscript{(22)}, quatre
\VUgras{rameurs} \textsuperscript{(24)} rament de toutes leurs forces, tandis
que le \VUgras{barreur} \textsuperscript{(26)} tient le \VUgras{gouvernail}
\textsuperscript{(25)} et dirige la frêle embarcation. Les \VUgras{rames}
\textsuperscript{(23)} frappent l'eau en cadence et les légères embarcations
filent à une vitesse vertigineuse.

%%K t03-c2-09-1 p
9. --- Quelques jeunes gens s'affrontent, près de la pile du pont, pour le prix
du \VUgras{mât de beaupré} \textsuperscript{(28)}. L'un d'eux avance
prudemment sur le mât souple et glissant pour attraper le petit
\VUgras{drapeau} \textsuperscript{(29)} fixé au bout. Son
\VUgras{concurrent} \textsuperscript{(30)} malchanceux est tombé à l'eau. Il
nage pour regagner la terre.

%%K t03-c2-10-1 p
10. --- Des jeunes gens plus âgés font des \VUgras{joutes nautiques}
\textsuperscript{(32)} près du \VUgras{quai} \textsuperscript{(33)}. Un
nombreux \VUgras{public} \textsuperscript{(31)} assiste à tous ces jeux, appuyé
sur le \VUgras{parapet} du \VUgras{pont} \textsuperscript{(27)} et massé sur la
\VUgras{rive}. Quelques spectateurs, pourtant, se retournent pour suivre le jeu
du \VUgras{mât de cocagne} \textsuperscript{(45)} ou pour admirer le
\VUgras{cortège du maire}, qui vient pour la \VUgras{remise} des prix.

%%K t03-c2-11-1 p
11. --- D'abord, voici quelques \VUgras{gamins} \textsuperscript{(35)} qui
précèdent les \VUgras{musiciens} \textsuperscript{(36)}; puis viennent le
\VUgras{maire} \textsuperscript{(37)}, les adjoints et le \VUgras{conseil
municipal} de la petite ville. Les \VUgras{pompiers} \textsuperscript{(38)}
marchent en dernier.

%%K t03-tit-12 sub
\VUsaut{5.0mm}
\VUpk{10.2pt}{40}{La foire.}
\VUsaut{3.6mm}

%%K t03-c2-12-1 p
12. --- Il y a beaucoup de monde sur le \VUgras{champ de foire}
\textsuperscript{(39)}. On vient voir les différentes \VUgras{baraques} : les
\VUgras{chevaux de bois} \textsuperscript{(40)}, le \VUgras{cirque}
\textsuperscript{(41)}, où la représentation a déjà commencé, le \VUgras{bal
public} \textsuperscript{(42)}, où les jeunes gens et les jeunes filles de la
ville goûtent le plaisir de la \VUgras{danse}.

%%K t03-c2-13-1 p
13. --- Dans le fond, on aperçoit les \VUgras{arènes} \textsuperscript{(44)},
où le vaillant \VUgras{matador} espagnol combat le \VUgras{taureau}
\textsuperscript{(45)} furieux, puis les \VUgras{montagnes russes}
\textsuperscript{(49)} et le \VUgras{stand de tir} \textsuperscript{(46)}, où
l'on s'exerce à manier les \VUgras{armes} et à tirer sur la \VUgras{cible}.

%%K t03-c2-14-1 p
14. --- Tout près, voici le \VUgras{théâtre forain} \textsuperscript{(47)}, où
l'on joue la \VUgras{comédie} et le \VUgras{drame}. Juste à côté se trouve la
baraque du \VUgras{géant} \textsuperscript{(48)} et du \VUgras{nain}. La
\VUgras{taille} du premier est de deux mètres quinze; le second est à peine
aussi grand qu'un enfant.

%%K t03-c2-15-1 p
15. --- Enfin, voici la grande \VUgras{ménagerie} \textsuperscript{(50)} avec
ses \VUgras{animaux sauvages}, son \VUgras{tigre} \textsuperscript{(51)} aux
\VUgras{griffes} puissantes, son \VUgras{lion} \textsuperscript{(52)} à
l'épaisse \VUgras{crinière}, ses deux \VUgras{girafes} \textsuperscript{(53)}
et son \VUgras{éléphant} \textsuperscript{(54)}, qui se sert si adroitement de
sa \VUgras{trompe}.

%%K t03-c2-16-1 p
16. --- Le \VUgras{dompteur} \textsuperscript{(55)} qui dresse ces animaux est
à la porte de la baraque.
\VUsaut{6.0mm}
\VUfilet{22mm}
